\section{Discussion}\label{sec:discussion}

\subsection{What the pilot establishes}

Three things, in decreasing order of confidence. The cost of screening a
record with a general-purpose model is between US\$0.0001 and US\$0.0016 at
current tariffs, which is small enough that the price of inference cannot be
the reason a review team declines to automate the first pass. The platform can
hold and audit a screening run of several thousand records with per-record
provenance, which the pilot exercised four times. And the agreement between
model and reviewer can be computed by the tool itself, which changes the
economics of validation: an agreement study stops being a separate research
project and becomes a button in the project dashboard.

\subsection{Error asymmetry is the design problem, not accuracy}

The pilot's most transferable observation is the shape of the errors, not their
rate. Across the two runs that produced interpretable matrices, false positives
outnumbered false negatives by roughly five to one. For a first-pass screen
this is the correct direction to err in, and it is also the reason why
reporting a single ``accuracy'' number for a screening tool is uninformative:
the same model can be tuned to miss nothing and to save little, or to save a
great deal and to miss studies. Sifter exposes the choice instead of concealing
it, through the three-way decision space and through the doubt queue, which
prioritises precisely the records on which the model declined to commit.

Whether the recall-oriented configuration is the right default for evidence
synthesis is an empirical question that this paper does not answer, and that
the field has not answered either. What the platform contributes is the
ability to state the configuration, measure its consequence and change it
without rewriting the tool.

\subsection{Provider independence as a methodological requirement}

Model behaviour changes with versions, prompts and sampling temperature, and
the guidance for generative AI submissions now requires the model version and
access date to be recorded \autocite{Farotimi_2025}. A tool that hard-codes one
provider cannot honour that requirement across a multi-year review, because
the model it used in the first year may not exist in the second. Sifter
therefore separates the protocol, the prompt template and the provider, and
stores the triplet with each decision. This is an engineering decision with a
methodological justification: it is what allows a review published in 2026 to
be re-screened in 2029 with the same protocol and a different model, and the
difference attributed to the model rather than to the protocol.

\subsection{Reporting against the guidance}

Table~\ref{tab:guidance} maps the journal's requirements for manuscripts
evaluating generative AI \autocite{Farotimi_2025} to the sections of this
paper. The last row is the honest one.

\begin{table}[htbp]
\centering
\caption{Reporting requirements for generative AI evaluations and where each is
addressed}
\label{tab:guidance}
\begin{tabular}{@{}p{0.34\textwidth}p{0.56\textwidth}@{}}
\toprule
Requirement & Where addressed \\
\midrule
Sampling and dataset characteristics &
Section~\ref{sec:evaluation}, Table~\ref{tab:corpus}; the 300-pair blind sample
is described in Section~\ref{sec:planned} \\
Prompt development and testing &
Section~\ref{sec:protocol}; prompts are versioned with the protocol \\
Established performance metrics &
Agreement, sensitivity and specificity in Table~\ref{tab:kappa}; WSS and NNR
outstanding \\
Benchmark against manual screening &
Table~\ref{tab:kappa}, with the provenance limitation stated \\
Quality assurance procedures &
Blind validation mode, Section~\ref{sec:blind} \\
Seeds and randomisation &
Sample selection is seeded and stored with the validation study \\
Computational environment and model versions &
Section~\ref{sec:system} and the per-decision record; versions and access dates
in the Data and Code Availability Statement \\
External validation &
Outstanding; Section~\ref{sec:planned} \\
\bottomrule
\end{tabular}
\end{table}

\subsection{Limitations}

The pilot has four limitations that a reader should carry into any use of the
numbers. It is a single installation used by a single reviewer, so
between-reviewer variability is not represented. The stored human decisions
were not made under a uniform protocol version and their timestamps were not
captured for every pair, which is why only one project yields an interpretable
$\kappa$. The evaluation covers title and abstract screening only; full-text
screening, risk of bias and data extraction are outside the scope of this
paper, although the platform has an extraction workspace. And the reference
standard in the agreement study is one human reviewer, not consensus between
two, which means the reported agreement is a lower bound on agreement with a
methodologically correct reference.

\subsection{Implications for practice}

For a review team, the practical implication is a change in what has to be
planned. The automated first pass does not remove the need for screening; it
changes the object that has to be validated. Instead of asserting that a tool
was used, the team can report how many records the tool screened, with which
model and prompt, at what cost, and with what agreement against a blind human
sample. Review protocols and PRISMA flow diagrams already have slots for part
of this information; the rest would need a reporting convention, and the
guidance cited above is a first step towards one.

For tool builders, the implication is that the unit of contribution is not the
prompt but the audit trail. A prompt that screens well on a benchmark and
leaves no record of what it saw is not reusable in evidence synthesis, however
good its numbers are.
